\chapter{Semigroupoid Paths and the Invalid-Inference Analysis}
\label{ch06:main}

Chapter~\ref{ch05:main} developed the licensed side of iteration: a single
endomap on a single object, living in a genuine monoid, with powers and a
finite idempotence theorem. This chapter develops its deliberately
crippled sibling. A \emph{semigroupoid path} is a nonempty sequence of
arrows whose adjacent objects match; it has a typed composite --- the
category laws guarantee that composing along the sequence is well defined
and association-independent --- and it has, by construction, \emph{no power
operation}. There is no $p^{2}$, no $p^{n}$, no ``run the path again.'' The
omission is not an unimplemented feature; it is the central design
statement of the kernel, and the reason the two constructions exist as
distinct runtime types rather than as one type with a flag.

Why spend a chapter on an absence? Because the absence encodes a theorem-
shaped mistake that recurs, in the wild, wherever operator formalisms meet
informal reasoning about repetition. The mistake has a fixed anatomy: take
a \emph{sequence} of operations --- typed, heterogeneous, possibly
time-varying --- and reason about it as if it were the \emph{iteration of
one} operation, thereby importing conclusions (fixed points, idempotence,
convergence, periodicity) that are theorems only about single elements of
finite monoids. Chapter~\ref{ch05:main} proved those theorems in their
legitimate habitat; this chapter shows, in three extended case studies,
exactly how each step of the illegitimate transfer fails in the kernel ---
not rhetorically but mechanically, with the precise construction site at
which a typed error is raised, or the precise operation that does not
exist in the API.

The chapter also situates the path construction against its standard
mathematical background --- the free category on a directed graph, in the
sense of Mac Lane \cite{maclanecwm} --- and closes with the boundary
statement that paths carry no dynamics and no semantics: a path is a
bookkeeping object for typed composition, and nothing about workflows,
processes, or cognition follows from one.

\section{The Path Construction}
\label{ch06:sec:construction}

\subsection{Definition}

\begin{definition}[Semigroupoid path]
\label{ch06:def:path}
A \emph{semigroupoid path} is a nonempty finite sequence of arrows
\[
  X_0 \xrightarrow{\;f_1\;} X_1 \xrightarrow{\;f_2\;} X_2
  \longrightarrow \cdots \xrightarrow{\;f_n\;} X_n ,
  \qquad n \ge 1,
\]
such that for each adjacent pair, the declared target of $f_k$ equals the
declared source of $f_{k+1}$ \emph{as objects} --- name and ordered carrier
both, exactly the definedness condition of
Definition~\ref{ch04:def:comp}. The path's \emph{source} is $X_0$ and its
\emph{target} is $X_n$; the objects $X_1, \dots, X_{n-1}$ are its
\emph{intermediate objects}. The constructor validates every adjacency at
construction time and raises the typed composition error, naming the
offending segment, if any adjacency fails; it also rejects the empty
sequence and any unordered or one-shot container of arrows, per the
ordered-input discipline of Chapter~\ref{ch03:main}.
\end{definition}

\begin{definition}[Path composite]
\label{ch06:def:composite}
The \emph{composite} of the path in Definition~\ref{ch06:def:path} is the
arrow
\[
  f_n \circ f_{n-1} \circ \cdots \circ f_1 \;:\; X_0 \to X_n ,
\]
computed by left-to-right folding of the kernel's composition. By the
associativity clause of Theorem~\ref{ch04:thm:laws}, every bracketing of
the fold yields the same arrow, so the composite is well defined
independently of evaluation order.
\end{definition}

\begin{ruleenv}[No power operation]
\label{ch06:rule:nopower}
The path type exposes construction, source, target, the composite, and
appending a type-compatible arrow. It exposes no exponentiation of any
kind: no operation takes a path and a natural number and returns
``the path repeated.'' This is an API absence enforced by the type system
of the host language: the power operation is defined only on the endomap
wrapper (Definition~\ref{ch05:def:endomap}), the path type is not that
type and is not convertible to it, and the endomap constructor accepts only
a single arrow whose source and target are equal as objects.
\end{ruleenv}

\begin{remark}[Why ``semigroupoid'']
\label{ch06:rem:semigroupoid}
The name records the algebraic habitat precisely. A semigroupoid (also
called a semicategory) has objects and typed composition, associative
where defined, but is not required to have identities at each object, and
its composable sequences do not form monoids unless source and target
coincide. The kernel's ambient structure $\cC$ is a full category ---
identities exist --- but a \emph{path} is an element of the free
composable-sequence structure over $\cC$'s arrows, and the operations the
kernel grants to paths are exactly the semigroupoid operations: composition
along matching types, and nothing that presupposes an enclosing monoid.
Powers presuppose an enclosing monoid. Hence their absence.
\end{remark}

\subsection{Relation to free categories on graphs}
\label{ch06:sec:free}

The path construction is a disciplined instance of a completely standard
object. Given a directed graph $G$ --- vertices, edges, each edge with a
source and target vertex --- the \emph{free category} $\mathcal{F}(G)$ on
$G$ has the vertices as objects and, as morphisms, the finite composable
strings of edges, including one empty string at each vertex serving as its
identity; composition is concatenation \cite{maclanecwm}. The
construction is left adjoint to the forgetful functor from small
categories to directed graphs, and its universal property is the
practical content of that adjunction: for any category $\mathcal{D}$ and
any graph morphism $\phi$ from $G$ into the underlying graph of
$\mathcal{D}$ --- an assignment of objects to vertices and morphisms to
edges, respecting sources and targets --- there is a \emph{unique}
functor $\hat{\phi} : \mathcal{F}(G) \to \mathcal{D}$ extending $\phi$,
namely the one sending each composable edge string to the composite of
its images,
\[
  \hat{\phi}\bigl(e_1 e_2 \cdots e_n\bigr)
  = \phi(e_n) \circ \cdots \circ \phi(e_2) \circ \phi(e_1),
\]
and the empty string at a vertex $v$ to $\id_{\phi(v)}$. Freeness means
the strings satisfy no equations beyond associativity and identity
bookkeeping: distinct strings remain distinct morphisms of
$\mathcal{F}(G)$ even when their images under some $\hat{\phi}$
coincide --- which is the categorical statement of the forgetting
phenomenon exhibited concretely in Example~\ref{ch06:ex:forgetting}.

Two comparisons locate the kernel's paths in this picture.

First: the kernel's paths are, up to the excluded empty string, the
morphisms of the free category on the underlying graph of $\cC$ (vertices:
the finite objects; edges: the arrows). The adjacency condition of
Definition~\ref{ch06:def:path} is exactly composability of edge strings.
The kernel excludes the empty path deliberately: an empty path would need
a distinguished object to serve as its source and target --- data not
derivable from an empty arrow sequence --- and the identity arrows of
$\cC$ already serve every legitimate use of ``do nothing at $X$.''
Nonemptiness keeps the path type's source and target total functions of
its data, with no special case.

Second, and more importantly: the kernel's path is a \emph{formal string
that remembers its factorization}, whereas its composite is a single arrow
that forgets it. The evaluation functor from the free category to $\cC$
--- the unique functor extending the identity assignment on vertices and
edges, per the universal property \cite{maclanecwm} --- is exactly the
kernel's \texttt{composite} operation. Distinct paths can have equal
composites (the fold collapses the string), and the kernel keeps both
artifacts available because they answer different questions: the path
answers ``what steps, in what order, through what intermediate types,''
which is the auditability question; the composite answers ``what total
function resulted,'' which is the denotational question. Nothing is ever
inferred about the path from properties of its composite --- that
direction of inference is precisely where the invalid inferences of
Section~\ref{ch06:sec:cases} live.

\begin{notation}
\label{ch06:not:path}
For a path $p$ with arrows $f_1, \dots, f_n$ we write $\llbracket p
\rrbracket = f_n \circ \cdots \circ f_1$ for its composite, mirroring the
denotation brackets used for typed words in the normal-form chapter; both
are evaluation maps from syntax-remembering objects to arrows.
\end{notation}

\section{Path Algebra}
\label{ch06:sec:algebra}

Although paths have no powers, they are not algebraically inert: they
concatenate, they extend one arrow at a time, and their evaluation
respects both operations. This section proves the three small facts the
rest of the monograph uses, and exhibits the forgetting phenomenon ---
distinct paths with equal composites --- that motivates keeping paths and
composites as separate artifacts.

\subsection{Concatenation}

\begin{definition}[Concatenation]
\label{ch06:def:concat}
Let $p$ be a path with arrows $f_1, \dots, f_m$ from $X_0$ to $X_m$, and
$q$ a path with arrows $g_1, \dots, g_n$ from $Y_0$ to $Y_n$. The
\emph{concatenation} $p \cdot q$ is defined if and only if $X_m = Y_0$ as
objects, in which case it is the path with arrows
$f_1, \dots, f_m, g_1, \dots, g_n$ from $X_0$ to $Y_n$.
\end{definition}

Definedness of the concatenation is a single adjacency check --- the same
comparison as arrow composition, at the junction --- and the result is
well formed because all its other adjacencies are inherited from $p$ and
$q$. In the implementation, concatenation is available as iterated
\texttt{append}; the definition above is its batch form.

\begin{proposition}[Paths form a semigroupoid]
\label{ch06:prop:semicat}
Concatenation is associative wherever defined: if $p \cdot q$ and
$q \cdot r$ are defined, then so are $(p \cdot q) \cdot r$ and
$p \cdot (q \cdot r)$, and they are equal.
\end{proposition}

\begin{proof}
Definedness: the junctions of $(p \cdot q) \cdot r$ are those of
$p \cdot q$ (the $p$/$q$ junction) and the $q$/$r$ junction, since
concatenation changes neither the source nor the target of the $q$
segment; symmetrically for the other association. Equality: both
associations have the same underlying arrow sequence --- concatenation of
sequences is associative --- and a path is determined by its arrow
sequence (the objects are read off the arrows). Hence the two paths are
equal as values.
\end{proof}

\begin{proposition}[Evaluation respects concatenation]
\label{ch06:prop:evalconcat}
If $p \cdot q$ is defined, then
$\llbracket p \cdot q \rrbracket = \llbracket q \rrbracket \circ
\llbracket p \rrbracket$.
\end{proposition}

\begin{proof}
Write the arrows as in Definition~\ref{ch06:def:concat}. Both sides are
composites of the sequence $f_1, \dots, f_m, g_1, \dots, g_n$ in order:
the left side by the fold of Definition~\ref{ch06:def:composite} over the
whole sequence, the right side by folding each half and then composing the
results. By the associativity clause of Theorem~\ref{ch04:thm:laws},
every bracketing of a fixed composable sequence yields the same arrow
(formally: by induction on $n$, moving one arrow at a time across the
central composition using associativity), so the two sides agree. All
intermediate composites are defined because every adjacency in the
combined sequence holds.
\end{proof}

\begin{lemma}[Append form]
\label{ch06:lem:append}
For a path $p$ and an arrow $g$ with source equal to $p$'s target,
$\llbracket p \cdot g \rrbracket = g \circ \llbracket p \rrbracket$,
where $p \cdot g$ abbreviates appending $g$ to $p$.
\end{lemma}

\begin{proof}
The case $n = 1$ of Proposition~\ref{ch06:prop:evalconcat}, with $q$ the
one-arrow path $(g)$, whose composite is $g$ itself.
\end{proof}

In categorical language, Propositions~\ref{ch06:prop:semicat}
and~\ref{ch06:prop:evalconcat} say that nonempty paths form a
semicategory (identities excluded by the nonemptiness convention of
Section~\ref{ch06:sec:free}) and that evaluation
$p \mapsto \llbracket p \rrbracket$ is a semifunctor into $\cC$ --- the
restriction to nonempty strings of the evaluation functor from the free
category \cite{maclanecwm}. Nothing deeper is claimed; the propositions
exist so that downstream code may reassociate path-building operations
freely, knowing the denotation cannot shift.

\subsection{Forgetting: distinct paths, equal composites}

\begin{example}[Two factorizations of one function]
\label{ch06:ex:forgetting}
Let $X = (\texttt{pair}, (0, 1))$ and let
$Y = (\texttt{triple}, (a, b, c))$ and
$Z = (\texttt{single}, (\ast))$ be auxiliary objects. Consider two paths
from $X$ to $X$:
\begin{enumerate}
\item $p_1$: through $Y$, with $u_1 : X \to Y$ mapping $(0, 0)$ and
$d_1 : Y \to X$ mapping $(1, 0, 0)$; composite
$\llbracket p_1 \rrbracket$ has tuple
$\bigl(d_1(0), d_1(0)\bigr) = (1, 1)$.
\item $p_2$: through $Z$, with $u_2 : X \to Z$ mapping $(0, 0)$ and
$d_2 : Z \to X$ mapping $(1)$; composite $\llbracket p_2 \rrbracket$ has
tuple $(1, 1)$.
\end{enumerate}
The two paths are unequal as paths --- different lengths would already
suffice, but here even the lengths agree and the arrow sequences differ
--- while their composites are extensionally equal: both denote the
constant map at index $1$. Every question about ``what happened along the
way'' (which intermediate object, what fan-out at the middle stage)
distinguishes $p_1$ from $p_2$; no question about the resulting total
function does.
\end{example}

The example is elementary, but its consequence for system design is not:
any pipeline audit that stores only composites has irreversibly discarded
the factorization, and any tool that needs the factorization --- tracing,
staging analysis, per-step attribution --- must consume paths. The kernel
therefore treats ``compute the composite'' as a lossy, explicit,
caller-initiated operation, never an implicit normalization. This is the
same design posture as the name-inclusive object identity of
Chapter~\ref{ch03:main}: information declared by the caller is never
silently collapsed, and every collapse is an auditable call site.

\begin{remark}[Paths and the normal-form machinery]
\label{ch06:rem:nfpreview}
The typed words of the normal-form chapter are close cousins of paths ---
sequences of generator labels with matching types --- and the closure
algorithm there is, in these terms, a systematic exploration of all paths
over a declared generator catalog, quotiented by extensional equality of
composites and represented by shortlex-least witnesses. The forgetting
phenomenon of Example~\ref{ch06:ex:forgetting} is exactly what the
closure's denotation table records: many words, one denotation, one
canonical representative. The present chapter's machinery is the
untabulated substrate of that construction.
\end{remark}

\section{The Invalid Inference}
\label{ch06:sec:invalid}

We now name the mistake the separation exists to block, in its general
form, before dissecting three concrete instances.

\begin{finding}[The sequence-to-power fallacy]
\label{ch06:find:fallacy}
Let $g_1, g_2, \dots, g_n$ be a sequence of operations applied in temporal
order. The \emph{sequence-to-power fallacy} is any inference of the form:
``the system after $n$ steps is $g^{n}$ for some single operation $g$;
therefore conclusions about powers of a single finite endomap ---
eventual periodicity, the factorial idempotence of
Theorem~\ref{ch05:thm:factorial}, existence of a unique idempotent in the
generated subsemigroup (Proposition~\ref{ch05:prop:monogenic}) --- apply
to the $n$-step system.'' The inference requires, and the informal
notation silently assumes, three separate facts: (i) all $g_t$ act on one
and the same object; (ii) all $g_t$ are equal to a single $g$; and
(iii) that $g$ is an endomap. Each can fail independently, and the kernel
is engineered so that whichever one fails first does so as a construction-
time type error or as an operation that does not exist, rather than as a
false theorem silently applied.
\end{finding}

The taxonomy is worth fixing in tabular form before the case studies,
because the three conditions fail at three different kernel surfaces, and
keeping the correspondence straight is what makes the studies
transferable to new instances:

\begin{center}
\small
\begin{tabular}{@{}llll@{}}
\toprule
Condition assumed & Fails in study & Kernel surface that objects &
Failure form \\
\midrule
(i) one object throughout & (a) & composition / endomap certification &
typed error \\
(ii) one generator repeated & (b) & no root-extraction operation exists &
API absence \\
(iii) a certified endomap & (c) & endomap constructor; no path power &
typed error / absence \\
\bottomrule
\end{tabular}
\end{center}

Note the two distinct failure forms. Conditions that assert a false
\emph{equality of objects} fail loudly, as typed errors at a
construction site --- there is a place in the program where the
identification is attempted, and that place raises. The condition that
asserts a false \emph{existence} (of a single generator $g$ with the
right power) fails silently in a stronger sense: the kernel simply
provides no operation whose use would presuppose it, so the assumption
has nowhere to be written down. Both forms serve the same audit goal by
different means, and the case studies illustrate the difference in
texture.

The three case studies below realize the three failure points in order:
(a) fails condition (i) --- the state spaces change over time; (b) fails
condition (ii) --- the objects agree and the composite is even an endomap,
but no single generator reproduces it; (c) attempts to import
Theorem~\ref{ch05:thm:factorial} wholesale into a pathful setting and
founders on (iii) at the endomap certification boundary. In each study we
give the setup, the tempting informal claim, and then --- the load-bearing
part --- the exact kernel step at which the formalization of the claim
becomes unwritable, with the specific error and the specific construction
that raises it.

\section{Three Case Studies}
\label{ch06:sec:cases}

\subsection{Case study (a): a time-varying workflow across changing state
spaces}
\label{ch06:sec:case-a}

\begin{example}[Setup]
\label{ch06:ex:case-a}
A three-stage pipeline processes items whose representation changes at each
stage. Declare three objects:
\begin{align*}
  S_0 &= (\texttt{intake}, (\texttt{raw\_a}, \texttt{raw\_b},
          \texttt{raw\_c})), \\
  S_1 &= (\texttt{screened}, (\texttt{keep}, \texttt{drop})), \\
  S_2 &= (\texttt{archived}, (\texttt{bin\_0}, \texttt{bin\_1},
          \texttt{bin\_2})),
\end{align*}
and two processing arrows: $g_1 : S_0 \to S_1$ with mapping $(0, 0, 1)$
(items \texttt{raw\_a} and \texttt{raw\_b} are kept, \texttt{raw\_c}
dropped), and $g_2 : S_1 \to S_2$ with mapping $(2, 0)$ (kept items to
\texttt{bin\_2}, dropped items to \texttt{bin\_0}). The path
$p = (g_1, g_2)$ is well formed --- the single adjacency check compares
$S_1$ against $S_1$ and passes --- and its composite is
$\llbracket p \rrbracket = g_2 \circ g_1 : S_0 \to S_2$ with mapping
$(2, 2, 0)$.
\end{example}

\emph{The tempting claim.} In an untyped operator notation, one writes the
pipeline as a single operator $T$ and reasons: ``running the pipeline
repeatedly, the system must stabilize; indeed $T^{k}$ is eventually
idempotent by finiteness'' --- an appeal to exactly the mathematics of
Chapter~\ref{ch05:main}, with $T$ playing the role of $f$.

\emph{Where the kernel stops it.} The formalization must produce the object
on which $T$ iterates, and there is none. Concretely, each of the three
available moves fails at a stated site:

\begin{enumerate}
\item \emph{Iterate the path.} The expression ``$p^{2}$'' does not
typecheck against any operation: the path type has no power
(Rule~\ref{ch06:rule:nopower}). This failure is an absence, not an
exception --- the program cannot be written.
\item \emph{Certify the composite as an endomap.} The expression
$\mathrm{Endomap}(\llbracket p \rrbracket)$ raises the typed composition
error at the certification check of Definition~\ref{ch05:def:endomap}:
the composite has declared source $S_0 = \texttt{intake}$ and declared
target $S_2 = \texttt{archived}$, which are unequal as objects (they differ
in name \emph{and} carrier --- sizes $3$ and $3$, but different element
sequences; and even equal carriers would not save it, by
Rule~\ref{ch05:rule:samecarrier}). The error message carries both declared
types, so the audit trail shows exactly which identification the informal
argument assumed and never justified.
\item \emph{Compose the pipeline with itself.} The expression
$\llbracket p \rrbracket \circ \llbracket p \rrbracket$ --- or
equivalently, appending $g_1$ to the path $p$ to begin a second lap ---
raises the composition error of Invariant~\ref{ch04:inv:mismatch}: the
would-be middle objects are $S_2$ (target of the first lap) and $S_0$
(source of the second), and $S_2 \ne S_0$. The failed adjacency is the
precise formal location of the informal argument's gap: ``run it again''
presupposes an arrow $S_2 \to S_0$ returning archived items to the intake
representation, and no such arrow was ever declared.
\end{enumerate}

\emph{The repair, if intended.} If the modeler genuinely intends re-intake,
the kernel does not forbid it --- it demands that the intent be an artifact.
Declare a re-intake arrow $r : S_2 \to S_0$ (choosing, and thereby
documenting, what each bin becomes on re-entry); then
$r \circ \llbracket p \rrbracket$ is an endomap on $S_0$,
$\mathrm{Endomap}(r \circ \llbracket p \rrbracket)$ certifies, and
Theorem~\ref{ch05:thm:factorial} applies to \emph{that} arrow, with the
conclusion depending --- as it must --- on the declared content of $r$.
The fallacy was never ``iteration is impossible''; it was ``iteration for
free, with the return map left implicit and unexamined.''

\subsection{Case study (b): a heterogeneous composite that happens to be
endo, confused with a generator's power}
\label{ch06:sec:case-b}

The second failure point is subtler: the types all line up, the composite
genuinely is an endomap, and the fallacy is the identification of that
composite with a \emph{power of a single generator} --- an equality claim
between arrows, false in general, whose falsity the kernel makes checkable
rather than assumable.

\begin{example}[Setup]
\label{ch06:ex:case-b}
Declare $X = (\texttt{register}, (0, 1))$ and
$Y = (\texttt{buffer}, (\texttt{lo}, \texttt{mid}, \texttt{hi}))$, with
arrows $h_1 : X \to Y$, mapping $(2, 0)$ (register $0$ to \texttt{hi},
register $1$ to \texttt{lo}), and $h_2 : Y \to X$, mapping $(0, 1, 1)$
(\texttt{lo} to register $0$; \texttt{mid} and \texttt{hi} to register
$1$). The path $q = (h_1, h_2)$ is well formed, and its composite
$c = \llbracket q \rrbracket = h_2 \circ h_1 : X \to X$ computes to
\[
  c(0) = h_2\bigl(h_1(0)\bigr) = h_2(2) = 1,
  \qquad
  c(1) = h_2\bigl(h_1(1)\bigr) = h_2(0) = 0 :
\]
the composite is the \emph{transposition} $\tau : X \to X$ with mapping
$(1, 0)$. Since source and target are both $X$,
$\mathrm{Endomap}(\tau)$ certifies without complaint. Everything so far
is legitimate.
\end{example}

\emph{The tempting claim.} ``The round trip through the buffer is a
two-step process; so the system state after the round trip is $g^{2}$ for
the underlying step operation $g$, and therefore'' --- here the fallacy
cashes out --- ``it inherits the structural properties of squares: for
instance, its cycle structure is constrained to that of a square of an
endomap on two points.''

\emph{Why the claim is false, concretely.} On the two-point carrier there
are exactly $2^{2} = 4$ endomaps: the identity $\id$ (mapping $(0,1)$),
the transposition $\tau$ (mapping $(1,0)$), and the two constants $c_0$
(mapping $(0,0)$) and $c_1$ (mapping $(1,1)$). Their squares are
\[
  \id^{2} = \id, \qquad
  \tau^{2} = \id, \qquad
  c_0^{2} = c_0, \qquad
  c_1^{2} = c_1 .
\]
The set of squares is therefore $\{\id, c_0, c_1\}$, and $\tau$ --- the
actual round-trip composite of Example~\ref{ch06:ex:case-b} --- \emph{is
not the square of any endomap on $X$}. The two-step round trip through a
three-point buffer realizes a map that no doubled one-step endo-process on
the register can realize. The informal claim is not merely unproved; it is
refuted by a four-case enumeration.

\emph{Where the kernel stops it.} The fallacy here is an arrow equality
--- $c = g^{2}$ for some endomap $g$ --- and equalities between arrows are
exactly what the kernel makes decidable and cheap: arrow equality is
extensional tuple comparison (Definition~\ref{ch04:def:extensional}). The
claim ``$c$ is a square'' becomes a finite search over the $|X|^{|X|}$
candidate generators, each tested by one composition and one comparison,
and it returns false. What the kernel refuses to supply is any construct
that would let the claim be \emph{assumed}: the path $q$ is not an endomap
and has no power; the certified endomap $\mathrm{Endomap}(c)$ has powers,
but its powers are powers \emph{of $c$}, the audited composite --- the
API affords $c^{2}$, and nothing affords ``the $g$ such that
$g^{2} = c$'' because no such $g$ need exist. In the type-theoretic
reading: the kernel provides introduction forms for powers (build a
certified endomap, then exponentiate) but no elimination form that
extracts roots, since root extraction is a nontrivial, often unsatisfiable
search problem, not a structural operation.

\emph{The general lesson.} A factorization through foreign objects can
realize endomaps unreachable as powers of any single same-object
generator. The composite operation deliberately \emph{loses} the
factorization (Section~\ref{ch06:sec:free}); the path retains it; and no
kernel operation manufactures a fictitious single-generator history for a
composite. When downstream analysis needs the history, it must consume the
path, which is the artifact that actually has one.

\subsection{Case study (c): importing the idempotence theorem into a
pathful setting}
\label{ch06:sec:case-c}

The third study takes the fallacy at its most seductive: a direct citation
of Theorem~\ref{ch05:thm:factorial} against a pipeline, with the carrier
sizes even agreeing, so that every cardinality-based sanity check passes
and only the typed discipline objects.

\begin{example}[Setup]
\label{ch06:ex:case-c}
A processing loop alternates between two representations of a
three-element state space that are, by design, carried by \emph{equal}
ordered carriers under different declared names:
\begin{align*}
  W   &= (\texttt{working}, (s_0, s_1, s_2)), \\
  W'  &= (\texttt{checkpoint}, (s_0, s_1, s_2)),
\end{align*}
with a checkpoint arrow $u : W \to W'$, mapping $(0, 2, 2)$, and a resume
arrow $v : W' \to W$, mapping $(1, 1, 0)$. The four-arrow path
$m = (u, v, u, v)$ is well formed (adjacencies alternate $W'$ against
$W'$ and $W$ against $W$), and its composite
$\llbracket m \rrbracket : W \to W$ is a certified endomap after wrapping.
So far, again, everything is legitimate --- and note that this loop, unlike
case (a), \emph{does} close.
\end{example}

\emph{The tempting claim.} ``Each pass of the loop is one step; by the
factorial idempotence theorem, after $3! = 6$ steps the process is
idempotent --- so the six-step behavior is a stable projection, and we may
reason about `the eventual state' of the pipeline.'' The citation then
gets applied not to $\llbracket m \rrbracket$ but to the \emph{stepwise}
description: the reasoner takes ``one step'' to mean ``one arrow of the
path'' and computes with $u$ and $v$ as if each were an iterate of a
single $f$.

\emph{Where the kernel stops it --- three distinct sites.} Formalizing the
claim forces a choice of which arrow the theorem is about, and every
choice that matches the informal wording fails at a typed boundary:

\begin{enumerate}
\item \emph{The theorem applied to $u$.} $\mathrm{Endomap}(u)$ raises the
typed error: $u$ has declared source \texttt{working} and declared target
\texttt{checkpoint}, unequal as objects \emph{despite equal ordered
carriers} --- this is Rule~\ref{ch05:rule:samecarrier}, and it is the
exact refusal that a carrier-only or cardinality-only identity discipline
would have waved through. Under such a laxer discipline, $u$ would certify
as an ``endomap'' on the shared carrier, $u^{6}$ would compute, and the
resulting arrow would be an artifact of an identification nobody declared:
the checkpoint-indexed values would be silently re-read as working-indexed
values on every second application. The kernel's name-inclusive identity
makes this precise failure unconstructible. The same applies
symmetrically to $v$.
\item \emph{The theorem applied to the path.} ``After six steps'' invites
$m^{?}$ or a six-arrow truncation reasoning; the path has no power
(Rule~\ref{ch06:rule:nopower}), and building the six-step object
explicitly means writing the six arrows down, at which point the
alternation --- $u, v, u, v, u, v$ --- is visible and the question ``six
steps of \emph{what single map}?'' answers itself: of none.
\item \emph{The theorem applied to the composite.} This is the one
legitimate reading, and the kernel supports it exactly:
$e = \mathrm{Endomap}(\llbracket m \rrbracket)$ certifies (the composite
is $W \to W$), the safe power $e^{3!} = e^{6}$ computes, and
Theorem~\ref{ch05:thm:factorial} guarantees its idempotence. But now the
claim has changed meaning, and the change is the entire point: the theorem
speaks about iterating \emph{the whole closed loop} $\llbracket m
\rrbracket$, i.e.\ about the $24$-arrow process ``four-arrow loop, six
times,'' not about six arrows of the original path, and not about any
property of $u$ or $v$ singly. The idempotence exponent is governed by the
carrier of the \emph{composite's} object ($|W| = 3$, so $3! = 6$ loop
iterations), a quantity with no relation to the path's length. The kernel
forces the reasoner to land on this reading because it is the only one
that constructs.
\end{enumerate}

\emph{The residue.} What survives formalization is a true, precise, and
noticeably weaker statement than the informal one: the closed loop's
composite, iterated $6$ times, is idempotent as a map on \texttt{working}.
What died in formalization is everything the informal statement smuggled:
that individual passes are iterates of one operation, that the number
``six steps'' refers to path arrows, and that stabilization happens
\emph{within} a single traversal. The kernel's contribution is that the
smuggling has to clear customs: each assumption either becomes a declared
arrow and a certified endomap, or it becomes a typed error with a named
construction site.

\section{Type-Level Enforcement}
\label{ch06:sec:enforcement}

The case studies leaned on two enforcement mechanisms; this section states
them as contracts, parallel to the mismatch-rejection contract of
Chapter~\ref{ch04:main}.

\begin{invariant}[Distinct runtime types, no coercion]
\label{ch06:inv:distinct}
The endomap wrapper and the path are distinct runtime types with disjoint
capabilities: the endomap type carries \texttt{power} and the idempotence
predicate and wraps exactly one arrow with certified equal source and
target; the path type carries adjacency-checked construction,
\texttt{composite}, and \texttt{append}, and wraps a nonempty arrow
sequence. Neither type is a subtype of the other; no implicit conversion
exists in either direction; and each constructor validates its own
invariant and raises on violation:
\begin{enumerate}
\item passing a path (or any non-arrow value) to the endomap constructor
fails its arrow-type validation;
\item passing an arrow with unequal declared source and target to the
endomap constructor raises the typed composition error naming the arrow
and its declared types;
\item passing an arrow sequence with any failed adjacency to the path
constructor raises the typed composition error naming the offending
segment;
\item the sanctioned bridge between the types is explicit and one-way:
take the path's \texttt{composite} (an ordinary arrow), then certify it
with the endomap constructor, which re-checks source-target equality from
scratch. There is no bridge that bypasses the check.
\end{enumerate}
\end{invariant}

\begin{checkclass}[Negative fixtures for the separation]
\label{ch06:check:separation}
The executable suite pins each clause of
Invariant~\ref{ch06:inv:distinct}: construction of an endomap from a
non-endo arrow (expected: typed composition error); construction of an
endomap from a source-target pair with equal carriers and distinct names
(expected: typed composition error --- the case-study (c) fixture);
construction of a path with one broken adjacency (expected: typed
composition error identifying the segment); construction of a path from an
empty sequence and from unordered containers (expected: validation error);
and the positive control, certifying a closed path's composite and taking
its factorial power (expected: success, with the idempotence assertion of
Chapter~\ref{ch05:main} holding). The positive control matters: a
separation that blocked the legitimate bridge as well as the illegitimate
shortcuts would be mere obstruction, and the suite demonstrates that the
sanctioned route stays open.
\end{checkclass}

\begin{remark}[Why not one type with a flag]
\label{ch06:rem:noflag}
A recurring design alternative deserves an explicit rejection note: a
single ``process'' type carrying a sequence of arrows and a boolean
\texttt{is\_closed}, with \texttt{power} raising at runtime when the flag
is false. This is behaviorally similar and structurally worse. The flagged
design makes ``path with a power method that sometimes works'' the
ambient shape, so every call site inherits a latent partiality, and the
error surfaces at the \emph{use} site (possibly deep in a computation,
possibly only on certain inputs) rather than at the \emph{construction}
site where the modeling decision was made. The two-type design moves the
failure to the earliest possible moment, makes the capability difference
legible in signatures --- a function that requires iterability asks for an
endomap in its type, and cannot be handed a path at all --- and lets the
host language's checker do statically what the flag design defers to
runtime. The kernel consistently prefers unrepresentable states to
checked-at-use states, and this is the flagship instance.
\end{remark}

\subsection{Stronger and weaker enforcement points on the design spectrum}
\label{ch06:sec:spectrum}

The two-type design sits at a particular point on a spectrum of
enforcement strength, and honesty about the spectrum clarifies both what
the kernel achieves and what it delegates.

\emph{Stronger: statically indexed types.} In a host language with
dependent or parametrically indexed types, the arrow type itself can carry
its source and target as type indices --- an arrow is a value of type
$\mathrm{Arrow}\,X\,Y$ --- so that an ill-typed composition or an endomap
claim about a non-endo arrow is rejected \emph{at compile time}, before
any program runs. The Lean model of the kernel lives at exactly this
strength: its composition operation demands a definitional match of the
middle object, and the invalid inferences of
Section~\ref{ch06:sec:cases} are not runtime errors there but terms that
fail to elaborate. The Python implementation cannot reach this point ---
the host's type system does not track value-level indices --- so it
realizes the discipline at the next point down.

\emph{The kernel's point: certified wrappers with construction-time
checks.} Every constructor validates its invariant once, eagerly, and the
resulting value is immutable, so possession of a value of the type is
possession of a checked certificate. The residual gap versus static
indexing is timing (violations surface when the offending construction
executes, not when it is written) --- the gap is not soundness: no
sequence of kernel API calls yields an uncertified endomap or an
unvalidated path, because the constructors are the only way in and the
values are frozen.

\emph{Weaker, rejected: conventions and documentation.} One step further
down, the discipline lives in docstrings (``callers should not iterate
heterogeneous pipelines'') and code review. The failure data motivating
this kernel --- the red-team findings summarized in
Chapter~\ref{ch01:main} --- came from systems at this level, where the
sequence-to-power fallacy is not blocked but merely deprecated. The
kernel's position is that a discipline worth stating is worth enforcing
at the strongest point the host allows, and that the distance between
``deprecated'' and ``unconstructible'' is where silent wrongness lives.

The ledger records the enforcement claim at its actual strength:
obligation PO-TKS8-CORE-03 (``keep one-carrier powers distinct from
heterogeneous paths'') carries status \textsc{executable validation} ---
separate runtime types plus the negative fixtures of
Checkclass~\ref{ch06:check:separation} --- not a formal-verification
status, because the statement ``no API sequence defeats the separation''
is an informal argument about the module's construction surface, checked
by inspection and by tests, not by proof. The Lean model's static
enforcement is verified in Lean's terms, but transferring that assurance
to the Python artifact is part of the open equivalence obligation
discussed in Chapter~\ref{ch04:main}.

\section{Boundary: What Paths Do Not Claim}
\label{ch06:sec:boundary}

\begin{boundary}
\label{ch06:bound:nodynamics}
A semigroupoid path is a bookkeeping object: a validated, ordered record of
typed arrows with a computable composite. It carries \emph{no} dynamics
claim --- nothing about time, causation, convergence, stability, or the
behavior of any process the path might be used to describe --- and
\emph{no} semantics claim: the kernel assigns no CAS, Noetic, cognitive,
or workflow meaning to any object, arrow, or path, and the ``workflow''
and ``pipeline'' vocabulary of this chapter's case studies is expository
narration about hypothetical modelers, not an asserted interpretation.
Mapping any canonical TKS/CAS construct onto kernel arrows or paths
remains an open, unlicensed obligation (ledger PO-TKS8-CORE-08,
\textsc{open / not licensed}); nothing in this chapter advances it. What
this chapter establishes is strictly structural: which compositions exist,
which iterations are licensed, and precisely where the unlicensed
inferences fail to construct.
\end{boundary}

\begin{boundary}
\label{ch06:bound:negative-scope}
The case-study analysis establishes that the three formalization attempts
fail at the stated kernel boundaries, as pinned by the negative fixtures of
Checkclass~\ref{ch06:check:separation}. It does not establish that the
kernel blocks every conceivable encoding of the sequence-to-power fallacy:
a determined user can always model around a type system --- for example,
by coercing everything into one giant declared object, at which point the
kernel's theorems apply honestly to that declared model and the modeling
choice itself is the auditable artifact. The claim is not that fallacies
are impossible; it is that this fallacy cannot be committed \emph{silently}
within the kernel's term language: every route to it requires constructing
an explicit identification (a return arrow, a declared shared object) that
then appears in the audit trail.
\end{boundary}

\section{Exercises}
\label{ch06:sec:exercises}

\begin{exercise}
\label{ch06:ex:adjacency}
For a path of $n$ arrows over objects of size at most $s$, count the
object-equality comparisons the constructor performs and the tuple
operations the composite performs, and conclude that path validation is
$O(n)$ object comparisons and evaluation is $O(n s)$ index operations.
Explain why validation cannot be deferred to evaluation time without
weakening the error contract (hint: a path may be constructed, stored, and
audited without ever being evaluated).
\end{exercise}

\begin{exercise}
\label{ch06:ex:closedcount}
Call a path \emph{closed} if its source equals its target. Over the
single object $X$ with $|X| = 2$ and the full arrow set
$\Hom(X, X)$ as available arrows, count the closed paths of length
exactly $3$ and the distinct composites they denote, and verify that the
denotation count is bounded by $|X|^{|X|} = 4$ while the path count is
$4^{3} = 64$ --- the forgetting ratio of
Example~\ref{ch06:ex:forgetting} made quantitative.
\end{exercise}

\begin{exercise}
\label{ch06:ex:square-roots}
Complete the enumeration sketched in Section~\ref{ch06:sec:case-b}:
for $|X| = 2$, verify by the four-case table that the squares in
$\Hom(X, X)$ are exactly $\{\id, c_0, c_1\}$, and then show that
\emph{every} endomap on two points arises as the composite of some
two-arrow path through a three-point object. Conclude that on two points,
factorization through a larger foreign object is strictly more expressive
than squaring, which is the mathematical core of case study (b).
\end{exercise}

\begin{exercise}
\label{ch06:ex:repair}
In case study (a), suppose the modeler declares the re-intake arrow
$r : S_2 \to S_0$ with mapping $(0, 1, 2)$ --- bins map back to raw items
positionally. Compute the endomap
$e = r \circ \llbracket p \rrbracket$ on $S_0$, analyze its functional
graph (tails, cycles, index, period) using the tools of
Chapter~\ref{ch05:main}, and find its minimal idempotent power. Then
change $r$'s mapping to $(2, 2, 2)$ and observe how the conclusion
changes --- demonstrating concretely that the iterated behavior depends
on the declared content of the return map, which is exactly the datum the
informal argument left implicit.
\end{exercise}

\begin{exercise}
\label{ch06:ex:emptypath}
The kernel excludes empty paths; the free category $\mathcal{F}(G)$
includes an empty path at each vertex \cite{maclanecwm}. Design the API
of an empty-path-admitting variant: what extra datum must the constructor
take, which operations change type, and where does
Proposition~\ref{ch06:prop:evalconcat} pick up side conditions? Weigh the
gained uniformity (paths form a category, not a semicategory) against the
cost (a partial function \texttt{source} becomes entangled with the extra
datum), and compare your conclusion with the kernel's choice in
Section~\ref{ch06:sec:free}.
\end{exercise}

\section{Summary}
\label{ch06:sec:summary}

\begin{center}
\small
\begin{tabular}{@{}p{5.6cm}p{7.2cm}@{}}
\toprule
Item & Status \\
\midrule
Path construction and composite
(Definitions~\ref{ch06:def:path}, \ref{ch06:def:composite}) &
Well-definedness from the category laws
(Theorem~\ref{ch04:thm:laws}); constructor-validated adjacencies;
executable validation. \\
Relation to free categories \cite{maclanecwm} &
Standard background; paths are the nonempty morphisms of the free category
on $\cC$'s underlying graph; composite is the evaluation functor. \\
Absence of path powers (Rule~\ref{ch06:rule:nopower}) &
API absence, enforced by distinct runtime types
(Invariant~\ref{ch06:inv:distinct}); negative fixtures
(Checkclass~\ref{ch06:check:separation}); ledger PO-TKS8-CORE-03,
\textsc{executable validation}. \\
Case studies (a)--(c)
(Sections~\ref{ch06:sec:case-a}--\ref{ch06:sec:case-c}) &
Each ill-typed step identified with its exact kernel error site; the
two-point square enumeration in (b) is a complete four-case check. \\
Dynamics and semantics &
None claimed (Boundary~\ref{ch06:bound:nodynamics}); semantic transfer
remains \textsc{open / not licensed} (PO-TKS8-CORE-08). \\
\bottomrule
\end{tabular}
\end{center}

With the composition layer (Chapter~\ref{ch04:main}), the licensed
iteration theory (Chapter~\ref{ch05:main}), and the unlicensed side's
enforcement (this chapter) in place, the kernel's finite-function core is
complete. The next chapter changes register entirely: it embeds the arrows
into finite-dimensional Hilbert spaces via the deterministic standard-basis
lift and proves exact operator-norm structure for the result.
